% Editorial draft — Results. Based on the existing manuscript's reported
% estimates, not a fresh computation from raw data. Historical runs only.
% Retains existing figure/table labels for later integration.
\section{Results}

\subsection{Final answers: heterogeneous models and systematic prompt shifts}

The four API models differ substantially in observed answers. Across the reported $T=0.01$ analytical sample, Hassner endorsement rates range from 10.3\% (Minimax) to 67.5\% (GLM-4.6), while Savary mean willingness ranges from 3.33 (Kimi) to 3.73 (Qwen). The Hassner treatment-by-model specification also improves model fit (McFadden pseudo-$R^2$: $0.507$ to $0.543$; $\Delta\mathrm{AIC}=-304$). Model identity therefore matters for both response levels and sensitivity to vignette manipulations. These descriptive rates pool different available stimulus conditions and should not be read as balanced model comparisons.

In the no-stimulus Hassner condition, support for sleep deprivation decreases from 44.6\% in the control to 19.3\% with protracted duration, 12.2\% with modest intelligence value, and 1.0\% with both modifications. These directions agree with the human experimental benchmark. In Savary, by contrast, the LLM mean rises from 3.39 in the private condition to 3.73 with public recognition and 4.09 with added social benefit. The increment associated with social benefit is directionally consistent with the source experiment, but the private-to-public increase reverses the human finding: public recognition \emph{reduced} willingness in the original study. The benchmark therefore reveals domain-specific limits of directional replication rather than a general success or failure of synthetic respondents (Figure~\ref{fig:replication}).

All eight historical system-prompt stimuli are associated with shifts \emph{away from the normatively approved response direction} in the pooled descriptive comparisons. In Hassner, the no-stimulus endorsement rate is 19.1\%; it is 41.6\% under C, 72.1\% under A, and ranges from 56.4\% to 90.0\% across H--M. In Savary, mean willingness is 3.74 without a stimulus, 3.44 under C, 3.32 under A, and 3.14--3.35 across H--M (Table~\ref{tab:incentive-means}; Figure~\ref{fig:stimuli-effects}). These are \emph{socially desirable response patterns} and \emph{stimulus-induced response shifts}, not direct estimates of human-style SDB or its mitigation.

The comparisons require a coverage qualification: H--M were fielded only on Kimi and GLM-4.6, while the baseline and core conditions have broader model coverage. Thus the displayed pooled means are useful descriptive summaries, but differences involving H--M are not automatically like-for-like causal estimates. In the historical attribute regression restricted to the two-model subsample, positive-framed L/M prompts are associated with substantially higher Hassner endorsement than negative-framed H--K (reported OR $\approx15$, $p<0.001$) and a smaller decrease in Savary willingness (reported $\beta=-0.13$, $p=0.02$). Because the observed prompts do not form a complete matched factorial, these are associations between prompt packages, not isolated effects of framing valence. Nonsignificant recipient and beneficiary coefficients do not establish that those attributes are inert.

Treatment-by-stimulus contrasts also vary: the original analysis reports moderation in some conditions, without reversal of the sign of the original Hassner treatment contrasts. In Savary, the private-to-public contrast remains positive across the eight historical stimuli (reported range $+0.05$ to $+0.26$), so none reproduces the human reversal under the tested settings. This result identifies a limitation of the \emph{observed prompts}, not a general impossibility of eliciting self-signaling-related reasoning.

\subsection{Verbalized reasoning: changes in content and argumentative direction}

The annotation results indicate that stimuli change the text models produce while selecting answers. Explicit \texttt{overcorrection\_concern} appears in approximately 9\% of no-stimulus traces and 72--76\% under A/C (expert-validation $\kappa=0.64$, $n=27$). The prevalence of explicitly rational-deliberative style rises from 0.67 to 0.86, and acknowledgement of multiple defensible positions from 0.57 to 0.80. The latter style codes were not included in the small human-validation set and should be treated as exploratory. These changes describe generated text, not independently verified cognitive awareness.

The A--C comparison is especially informative in Hassner. Endorsement is 72.1\% under A versus 41.6\% under C, while several coded category-prevalence differences are small. Direction codes, however, differ markedly: signed presence of \texttt{ratio\_of\_goals\_and\_means} changes from $-0.01$ under C to $+0.42$ under A, \texttt{majority\_reliance} from $+0.06$ to $+0.41$, and \texttt{cruelty} from $-0.68$ to $-0.34$ (Table~\ref{tab:direction-AvsC}). The first of these codes was not covered by expert validation; the others have limited or category-specific reliability evidence. The pattern is consistent with a shift in the \emph{polarity of verbalized arguments}, but it does not show that these arguments mediate the stimulus effect.

In Savary, both A and C strongly increase annotation of personal costs and benefits, while their mean final answers remain comparatively close (3.32 and 3.44). Across both domains, the annotation vocabulary therefore responds to prompts, but the relation between coded text and the answer is domain-dependent.

A separability benchmark classifies stimulus-present versus no-stimulus traces from annotation features at reported cross-validated ROC-AUC $0.96$, compared with $0.68$ for low- versus high-temperature no-stimulus traces (Figure~\ref{fig:benchmark}). The original analysis also reports that the contrast survives exclusion of directly stimulus-linked codes. This demonstrates that the two classification tasks have different separability under the chosen features and samples; it does not by itself measure relative causal impact on internal reasoning.

The annotation instrument has uneven validation. In a 27-trace expert-coded subset, agreement is high for \texttt{cruelty} and \texttt{personal\_cost\_burden} (both $\kappa=1.00$ for causality coding) and moderate-to-substantial for \texttt{overcorrection\_concern} ($\kappa=0.64$), whereas several psychological categories have near-zero agreement (Table~\ref{tab:kappa}). Interpretations of subtle self-presentational categories must therefore remain tentative.

\subsection{Associations between reasoning annotations and answers}

Adding annotation features increases reported explanatory fit in Savary from $R^2=0.419$ to $0.482$ ($\Delta R^2=0.063$). Overcorrection concern is negatively associated with willingness ($\beta=-0.047$, $p=0.012$), while positively directed personal-cost and private-interest codes are associated with higher willingness (both reported $\beta=+0.046$). A gradient-boosting model also outperforms the linear benchmark in cross-validation (reported $R^2=0.558$ versus $0.475$). These patterns indicate predictive information in verbalized reasoning beyond the specified treatment, stimulus, and model controls, but do not demonstrate that the generated text is causally upstream of the answer.

For Hassner, the original non-tautological annotation specification reports a smaller increment in fit ($\Delta R^2=0.052$). Positive directions of majority-reliance and means--ends arguments are associated with endorsement, while overcorrection concern itself is not a significant predictor ($p=0.92$). Taken together, the two experiments suggest that the \emph{observable association} between reasoning categories and final responses varies by domain. The small expert-validation sample and co-generation of traces and answers limit stronger process claims.

\subsection{Exploratory hidden-state results}

The local Qwen3-8B analysis covers 98 runs (56 Hassner and 42 Savary) under a different stimulus inventory from the API study. In matched comparisons, the reported cumulative hidden-state trajectory is shorter under stimulus conditions than without a stimulus (3{,}370 versus 3{,}444 distance units; paired Wilcoxon $p<0.001$; Figure~\ref{fig:trajectory}). This is a geometric difference across layer representations, not a direct measure of computational effort or deliberation.

The between-/within-condition separation statistic peaks near layers 20--21 in the reported comparisons. Linear CKA remains above 0.94 across layers (reported L20 CKA $=0.994$), indicating broadly similar representational geometry despite detectable condition differences (Figure~\ref{fig:separation}). These summaries do not uniquely identify a translation mechanism or a semantic processing stage.

The local model also exhibits a representational--behavioral dissociation: stimulus-related hidden-state differences are detectable even where its Hassner responses are near-saturated (reported $P(\mathrm{Yes})\approx0.998$). A treatment-status probe performs poorly under leave-one-group-out validation despite higher standard cross-validation performance, cautioning against scenario-specific leakage. For Savary, a condition-residualized ridge probe reaches a reported leave-one-out $R^2=0.665$ at layer 27 (Figure~\ref{fig:decodability}). The distinct layer peaks are descriptive: they do not establish when a decision is causally formed. Small samples, high-dimensional activations, and the separate model architecture constrain generalization.

\subsection{Synthesis}

Three results are most robust as a descriptive account of the historical corpus. First, model responses vary substantially across model families and between the two social domains; the Savary public-recognition contrast reverses the human benchmark. Second, system-prompt stimuli consistently accompany shifts away from the normatively approved response direction in the reported descriptive comparisons, while their apparent magnitude depends on prompt condition and model coverage. Third, stimuli are associated with changes in verbalized reasoning and with measurable hidden-state differences in a separate local model, but neither track alone identifies the mechanism producing final answers.

The historical findings support the study of system prompts as consequential experimental inputs, not as neutral wrappers. They do \emph{not} establish that steering responses away from social approval makes LLMs more human-like, that the observed attribute associations isolate framing valence, or that verbalized and representational measures constitute a causal chain. The planned N0/negative/positive follow-up will be integrated only after its results are available.
